\section{Conclusion}
\label{sec:conclusion}

% Discussion, limitations and future directions


% In this work, we explored using a masked diffusion model for the generation of chess puzzles. We followed the work of \citeauthor{feng2025generatingcreativechesspuzzles} and trained a masked diffusion model first with supervised training and later improved its performance by reinforcement learning. Additionally, we condition the generative process on the themes of a puzzle, for instance fork or mate-in-two, which steers the generation toward spesific puzzles. By supervised training we managed to train a model, which generates puzzles puzzles around $0.18\%$ of the time. We improve our model using reinforcement learning to generate puzzles around $0.37\%$ of the time. We also manage to improve the diversity of the generated positions.
%
% We had four main contributions from this research. Firstly, we demonstrated the effectiveness of masked diffusion models for generation of structured outputs. Secondly, we showed that the puzzle generation can be conditioned on specific themes. Thirdly, we demonstrated that the performance of masked diffusion models can be improved with reinforcement learning, although the task at hand is very difficult. Lastly, we demonstrated that the performance of the generation is improved by an auxiliary task of best move generation. Looking back at the contributions, although we have not demonstrated the generalization of these contributions to other tasks, in chess puzzle generation, the masked diffusion model works well and can be improved in many ways.
%
% For reinforcement learning, we employ an algorithm called DDPO. For this, we define a Markov decision process that works 
%
% For DDPO, we derive easy expressions for the step-wise KL divergence and the step-wise entropy.
%
% During this research, we have encountered many challenges. 
%
% This research has many limitations. (Generalizability, authors knowledge on chess etc.)
%
% In the future research, many questions remain. Firstly, we were unable to find a configuration, where the recently published ESPO algorithm would have worked. ESPO has offered benefits compared to DDPO in terms of efficiency and may improve the RL pipeline in this problem as well if a good configuration is found. Secondly, demonstrating similar setups in other NLP problems, and showcasing generalization, would push the research boundary further. Lastly, experimenting with different types of puzzle generation frameworks, such as combining iterative algorithms and generative models in one chain-of-thought \cite{wei2022chainofthought} system could improve puzzle generation and allow for easier generalization to other problems.
%
% \begin{itemize}
%     \item \textbf{Effectiveness of masked diffusion models:} Summarize how well the masked diffusion model performed at generating valid, structured chess puzzles compared to the baselines.
%     \item \textbf{Conditioning on themes:} Discuss the success rate of guiding the puzzle generation using specific tactical themes and any trade-offs observed (e.g., lower theme adherence after RL).
%     \item \textbf{Reinforcement learning improvements and challenges:} Highlight the increase in puzzle rate and counter-intuitiveness achieved through RL, but also mention the limitations encountered, such as mode collapse during extended training.
%     \item \textbf{Impact of auxiliary tasks:} Conclude on how predicting the best move simultaneously with the position improved the model's underlying chess knowledge and generative capabilities.
%     \item \textbf{Future directions:} Suggest applying this framework to other NLP tasks requiring strict structured outputs (like HTML rendering or code generation) and exploring better diversity filters or reward functions to prevent mode collapse in RL.
% \end{itemize}






% In this work, we explored using a masked diffusion model for the generation of chess puzzles, building upon the framework proposed by \citeauthor{feng2025generatingcreativechesspuzzles}. We initially trained a masked diffusion model using supervised training and subsequently enhanced its performance through reinforcement learning. A key aspect of our approach was conditioning the generative process on specific puzzle themes, effectively steering the generation toward desired puzzle types. After supervised training, our model successfully generated valid puzzles approximately $0.18\%$ of the time, although the puzzle rate changed a lot for each theme. By applying reinforcement learning, we more than doubled this success rate to around $0.37\%$, while simultaneously improving the diversity of the generated positions.

% We presented four main contributions in this research. First, we demonstrated the effectiveness of masked diffusion models for generating strict, structured outputs. Second, we showed that chess puzzle generation can be accurately conditioned on specific tactical themes. Third, we established that the performance of masked diffusion models can be meaningfully improved through reinforcement learning, despite the inherent difficulty of the task. Finally, we demonstrated that integrating the auxiliary task of predicting the best move alongside the position enhances the generative capabilities of the model. While we focused specifically on chess puzzle generation and did not demonstrate generalization to other tasks in this study, our results indicate that masked diffusion models can be very capable and offer various avenues for improvement in structured generation tasks.

% For the reinforcement learning phase, we employed the Denoising Diffusion Policy Optimization algorithm, which is extremely principled and theoretically grounded. To facilitate this, we formulated a Markov decision process following \citeauthor{black2024trainingdiffusionmodelsreinforcement}, which is suitable for the generation process. The MDP was formulated step-wise, which contrasts with the some of the latest dLLM research \cite{ou2025principledrldiffusionllms}. The states of the MDP were the partially unmasked token sequences and the actions were the generation of a new partially unmasked sequence. The transitions were deterministic and the rewards zero except for the last step where a nonzero reward was given to good puzzles and illegal positions. Additionally, we derived tractable expressions for the step-wise KL divergence and step-wise entropy, which are easy to compute and allow comparisons to traditional LLMs. With these expressions, we attempted to prevent model collapse and unrealistic positions.

% Despite the results, we encountered several challenges and limitations during this research. For instance, while we hope our findings are broadly applicable, we have not yet demonstrated the generalization of these contributions to other natural language processing tasks. Additionally, the manual evaluation and selection of the generated puzzles were inherently constrained by the author's level of chess knowledge. A stronger player or puzzle composer might offer a different perspective specifically to evaluation of our final model. Lastly, the stability and model collapse problems in reinforcement learning posed challenges to the algorithm we used. Connected to this, the decrease in theme match rate during RL may be caused by our theme match criterion, which is by no means the only way of defining a match. Perhaps, some combinations of themes allow the model to generate puzzles more often causing the model to converge to generating positions with the easy theme. These limitations leave room for further experimentation and studies.

% Looking ahead, numerous questions and directions for future research remain. During our experiments, we were unable to find a configuration where the recently published ESPO algorithm succeeded. Since algorithms that efficiently approximate the importance ratio have shown some benefits over DDPO \cite{ou2025principledrldiffusionllms, wang2026d2improvingresoning}, discovering a viable configuration could further optimize the RL pipeline for this problem. Furthermore, applying similar setups to other NLP domains to showcase generalizability would significantly advance the research boundary. Lastly, experimenting with alternative puzzle generation frameworks, such as combining iterative algorithms and generative models within a chain-of-thought system \cite{wei2022chainofthought}, could not only improve chess puzzle generation but also pave the way for easier adaptation to problems in other domains.

In this work, we explored using a masked diffusion model for the generation of chess puzzles, building upon the framework proposed by \citeauthor{feng2025generatingcreativechesspuzzles}. We initially trained a masked diffusion model using supervised training and subsequently enhanced its performance through reinforcement learning. A key aspect of our approach was conditioning the generative process on specific puzzle themes, effectively steering the generation toward desired puzzle types. After supervised training, our model successfully generated valid puzzles approximately $0.18\%$ of the time, although the puzzle rate varied considerably across different themes. By applying reinforcement learning, we more than doubled this success rate to around $0.37\%$, while simultaneously improving the diversity of the generated positions.

We presented four main contributions in this research. First, we demonstrated the effectiveness of masked diffusion models for generating strict, structured outputs. Second, we showed that chess puzzle generation can be accurately conditioned on specific tactical themes. Third, we established that the performance of masked diffusion models can be meaningfully improved through reinforcement learning, despite the inherent difficulty of the task. Finally, we demonstrated that integrating the auxiliary task of predicting the best move alongside the position enhances the generative capabilities of the model. While we focused specifically on chess puzzle generation and did not demonstrate generalization to other tasks in this study, our results indicate that masked diffusion models are highly capable and offer various avenues for improvement in structured generation tasks.

For the reinforcement learning phase, we employed the Denoising Diffusion Policy Optimization algorithm, an approach that is highly principled and theoretically grounded. To adapt this algorithm for our task, we formulated a Markov decision process following the methodology of \citeauthor{black2024trainingdiffusionmodelsreinforcement}. Crucially, our MDP was formulated step-wise, which contrasts with some of the latest research in discrete large language models \cite{ou2025principledrldiffusionllms, rojas2025improvingreasoningdiffusionlanguage}. In this setup, the states correspond to the partially unmasked token sequences, while the actions define the generation of the subsequent partially unmasked sequence. The transitions between states are deterministic, and the rewards remain zero until the final step, at which point a non-zero reward is assigned based on the quality of the puzzle. To further stabilize the training process, we derived tractable expressions for the step-wise KL divergence and step-wise entropy. These expressions allowed a direct comparison to traditional autoregressive LLMs, and served as vital regularization tools to prevent model collapse and the generation of unrealistic positions.

Despite our results, we encountered several challenges and limitations during this research. For instance, while we hope our findings are broadly applicable, we have not yet demonstrated the generalization of these contributions to other natural language processing tasks. Additionally, the manual evaluation and selection of the generated puzzles were inherently constrained by the author's level of chess knowledge. A stronger chess player or puzzle composer might offer a different perspective specifically regarding the evaluation of our final model. Furthermore, stability issues and model collapse during reinforcement learning posed significant challenges to our chosen algorithm. Relatedly, the observed decrease in the theme match rate during RL may stem from our specific theme match criterion, which is by no means the only way to define a successful match. It is possible that certain theme combinations inherently allow the model to generate puzzles more easily, causing it to converge toward generating positions that favor these "easier" themes. These limitations highlight clear avenues for further experimentation and study.

Looking ahead, numerous questions and directions for future research remain. During our experiments, we were unable to find a configuration where the recently published ESPO algorithm succeeded. Since algorithms that efficiently approximate the importance ratio have shown benefits over DDPO \cite{ou2025principledrldiffusionllms, wang2026d2improvingresoning}, discovering a viable ESPO configuration could further optimize the RL pipeline for this problem. Furthermore, applying similar setups to other NLP domains to showcase generalizability would significantly advance the research boundary. Lastly, experimenting with alternative puzzle generation frameworks, such as combining iterative algorithms and generative models within a chain-of-thought system \cite{wei2022chainofthought}, could not only improve chess puzzle generation but also pave the way for easier adaptation to complex problems in other domains.
